\scanpage{77}
coeff. entiers). Comme on sait que ce sous-groupe est $\mathrm{Sl}(2,\mathbb Z)$ tout entier, et que (si $\Lambda_0=\operatorname{Im}(\mathbb Z\mathbin{\unsure{-|\to}}\Lambda)$) \deleted{\unsure{celui}} l'hom
\[
 \mathrm{Sl}(2,\mathbb Z)\longrightarrow
 \underset{\substack{\rotatebox{-90}{$\subset$}\\\mathrm{Sl}(2,\Lambda)\subset\mathrm{Gl}(2,\Lambda)}}{\mathrm{Sl}(2,\Lambda_0)}
 \quad\text{est surjectif}
\]
on voit que le sous-groupe de $\mathrm{\deleted{S}Gl}(2,\Lambda)$ engendré par les $\tau_i^{r}$ n'est autre que \deleted{$\mathrm{Sl}(2,\Lambda_0)$.} \unsure{/}
\[
 \underset{\substack{\rotatebox{-90}{$\subset$}\\\mathrm{Sl}(2,\Lambda)}}{\mathrm{Sl}(2,\Lambda_0)}
 =\left\{\begin{pmatrix}a&b\\c&d\end{pmatrix}\ \middle|\
 \begin{aligned}&a,b,c,d\in\Lambda_0=\operatorname{Im}(\mathbb Z\mathbin{-}\Lambda)\\&ad-bc=1\end{aligned}\right\},
\]
en faisant opérer celui-ci \deleted{à droite} sur l'ens des repères, via
\[
 \deleted{\unsure{(f}}\ (x,y)\longmapsto(ax+by,cx+dy)
\]
Ainsi, on a prouvé

\uline{Lemme} Soit $u$ un autom. de $P_{M,i}$, qui fixe \inserted{deux éléments $x,y\in S_M$ formant} une $i$-base $(x,y)$ de $M$ (i.e. $x\wedge y\in i\subset\overrightarrow{\unsure{\pi|\Pi}}_M=\det M^*$ ---), \deleted{Alors} supposons de plus que $u$ commute à l'antipodisme $-\mathrm{id}_M$. Alors $u$ fixe \deleted{les} tous les éléments de $\widetilde S_M$ qui sont de la forme $ax+by$, avec $a,b\in\Lambda_0=\operatorname{Im}(\mathbb Z\to\Lambda)$, et premiers \deleted{\unsure{ce}} entre eux (i.e. $a\Lambda+b\Lambda=\Lambda$ i.e. $\exists c,d\in\Lambda$ avec $ad-\unsure{cc|ce}=1$).
% archive sheet number: 76
